%%%%%%%%%%%%%%%%
%%% deut 15 %%%
%%%%%%%%%%%%%%%%

\Chapter{15}

\Verse{1}
{
	\hDtn{מִקֵּץ שֶׁבַע שָׁנִים תַּעֲשֶׂה שְׁמִטָּה׃}
}
{
	\eDtn{“At the end of seven years you shall make a remission.}
}
{
}

\Verse{2}
{
	\hDtn{וְזֶה דְּבַר הַשְּׁמִטָּה שָׁמוֹט כּל בַּעַל מַשֵּׁה יָדוֹ אֲשֶׁר יַשֶּׁה בְּרֵעֵהוּ לֹא יִגֹּשׂ אֶת רֵעֵהוּ וְאֶת אָחִיו כִּי קָרָא שְׁמִטָּה לַיהֹוָה׃}
}
{
	\eDtn{
		And this is the matter of the remission: every holder of a
		loan is to remit what he has lent his neighbor. He shall
		not demand it of his neighbor and his brother, because a
		remission for YHWH has been called.
	}
}
{
%	\footnote{
%		demand. This is the same root that is understood in its noun form to
%		mean taskmasters in the story of the enslavement in Egypt (Exod 3:7;
%		5:10). As a verb it means to require, demand, or exact something from
%		someone.
%	}
}

\Verse{3}
{
	\hDtn{אֶת הַנּכְרִי תִּגֹּשׂ וַאֲשֶׁר יִהְיֶה לְךָ אֶת אָחִיךָ תַּשְׁמֵט יָדֶךָ׃}
}
{
	\eDtn{
		You may demand it of a foreigner, but your hand shall
		remit whatever of yours is with your brother.
	}
}
{
}

\Verse{4}
{
	\hDtn{אֶפֶס כִּי לֹא יִהְיֶה בְּךָ אֶבְיוֹן כִּי בָרֵךְ יְבָרֶכְךָ יְהֹוָה בָּאָרֶץ אֲשֶׁר יְהֹוָה אֱלֹהֶיךָ נֹתֵן לְךָ נַחֲלָה לְרִשְׁתָּהּ׃}
}
{
	\eDtn{
		Nonetheless, there won’t be an indigent one among you,
		because YHWH will bless you in the land that YHWH, your
		God, is giving you as a legacy, to take possession of it—
	}
}
{
}

\Verse{5}
{
	\hDtn{רַק אִם שָׁמוֹעַ תִּשְׁמַע בְּקוֹל יְהֹוָה אֱלֹהֶיךָ לִשְׁמֹר לַעֲשׂוֹת אֶת כּל הַמִּצְוָה הַזֹּאת אֲשֶׁר אָנֹכִי מְצַוְּךָ הַיּוֹם׃}
}
{
	\eDtn{
		only if you’ll listen to the voice of YHWH, your God, to
		be watchful to do all of this commandment that I command
		you today.
	}
}
{
}

\Verse{6}
{
	\hDtn{כִּי יְהֹוָה אֱלֹהֶיךָ בֵּרַכְךָ כַּאֲשֶׁר דִּבֶּר לָךְ וְהַעֲבַטְתָּ גּוֹיִם רַבִּים וְאַתָּה לֹא תַעֲבֹט וּמָשַׁלְתָּ בְּגוֹיִם רַבִּים וּבְךָ לֹא יִמְשֹׁלוּ׃}
}
{
	\eDtn{
		When YHWH, your God, will have blessed you as He spoke to
		you, then you’ll lend to many nations, but you won’t
		borrow; and you’ll dominate many nations, but they won’t
		dominate you.
	}
}
{
}

\Verse{7}
{
	\hDtn{כִּי יִהְיֶה בְךָ אֶבְיוֹן מֵאַחַד אַחֶיךָ בְּאַחַד שְׁעָרֶיךָ בְּאַרְצְךָ אֲשֶׁר יְהֹוָה אֱלֹהֶיךָ נֹתֵן לָךְ לֹא תְאַמֵּץ אֶת לְבָבְךָ וְלֹא תִקְפֹּץ אֶת יָדְךָ מֵאָחִיךָ הָאֶבְיוֹן׃}
}
{
	\eDtn{
		“When there will be an indigent one among you from one of
		your brothers within one of your gates in your land that
		YHWH, your God, is giving you, you shall not fortify your
		heart and shall not shut your hand from your brother who
		is indigent.
	}
}
{
}

\Verse{8}
{
	\hDtn{כִּי פָתֹחַ תִּפְתַּח אֶת יָדְךָ לוֹ וְהַעֲבֵט תַּעֲבִיטֶנּוּ דֵּי מַחְסֹרוֹ אֲשֶׁר יֶחְסַר לוֹ׃}
}
{
	\eDtn{
		But you shall open your hand to him and shall lend to him,
		enough for his shortage that he has.
	}
}
{
}

\Verse{9}
{
	\hDtn{הִשָּׁמֶר לְךָ פֶּן יִהְיֶה דָבָר עִם לְבָבְךָ בְלִיַּעַל לֵאמֹר קָרְבָה שְׁנַת הַשֶּׁבַע שְׁנַת הַשְּׁמִטָּה וְרָעָה עֵינְךָ בְּאָחִיךָ הָאֶבְיוֹן וְלֹא תִתֵּן לוֹ וְקָרָא עָלֶיךָ אֶל יְהֹוָה וְהָיָה בְךָ חֵטְא׃}
}
{
	\eDtn{
		Watch yourself in case there will be something
		good-for-nothing in your heart, saying, ‘The seventh year,
		the remission year, is getting close,’ and your eye will
		be bad toward your brother who is indigent, and you won’t
		give to him; and he’ll call to YHWH about you, and it will
		be a sin in you.
	}
}
{
}

\Verse{10}
{
	\hDtn{נָתוֹן תִּתֵּן לוֹ וְלֹא יֵרַע לְבָבְךָ בְּתִתְּךָ לוֹ כִּי בִּגְלַל הַדָּבָר הַזֶּה יְבָרֶכְךָ יְהֹוָה אֱלֹהֶיךָ בְּכל מַעֲשֶׂךָ וּבְכֹל מִשְׁלַח יָדֶךָ׃}
}
{
	\eDtn{
		You shall give to him, and your heart shall not be bad
		when you’re giving to him, because on account of this
		thing YHWH, your God, will bless you in everything you do
		and in everything your hand has taken on.
	}
}
{
}

\Verse{11}
{
	\hDtn{כִּי לֹא יֶחְדַּל אֶבְיוֹן מִקֶּרֶב הָאָרֶץ עַל כֵּן אָנֹכִי מְצַוְּךָ לֵאמֹר פָּתֹחַ תִּפְתַּח אֶת יָדְךָ לְאָחִיךָ לַעֲנִיֶּךָ וּלְאֶבְיֹנְךָ בְּאַרְצֶךָ׃}
}
{
	\eDtn{
		Because there won’t stop being an indigent in the land. On
		account of this I command you, saying: you shall open your
		hand to your brother, to your poor, and to your indigent
		in your land.
	}
}
{
%	\footnote{
%		there won’t stop being an indigent in the land. This is generally
%		understood to mean that there will never cease to be poor people. And
%		the problem is that just a few verses earlier the text says the
%		opposite: “there won’t be an indigent one among you” (15:4). Ibn Ezra
%		and others have explained this apparent contradiction by noting that the
%		statement that “there won’t be an indigent one among you” goes on to say
%		“only if you’ll listen.... ” Since Israel will not, in fact, listen in
%		the future, therefore “there won’t stop being an indigent in the land.”
%		Ramban criticizes this view, saying that it is unimaginable that the
%		Torah would say that the Jewish people will never come to listen to the
%		Torah’s commandments. It seems to me that the problem arises because
%		everyone takes this verse to mean that there will never stop being
%		indigent people. But it simply says “there won’t stop.” I take that to
%		mean that poverty will not just come to a stop on its own one
%		day—without any action by humans. That is one of the ways in which this
%		word is used elsewhere in the Torah. Sarah’s menopause is described as
%		“Sarah had stopped having the way of women” (Gen 18:11). The end of the
%		plague of hail in Egypt is described as “the shower and the hail
%		stopped” (Exod 9:34). So here, too, the full verse says that poverty
%		will not simply stop—the poor will not just disappear—and “On account of
%		this I command you, saying: you shall open your hand to your brother, to
%		your poor, and to your indigent in your land.” The statements here and
%		in v. 4 are consistent: there will be no poverty only if people act to
%		end it.
%	}
}

\Verse{12}
{
	\hDtn{כִּי יִמָּכֵר לְךָ אָחִיךָ הָעִבְרִי אוֹ הָעִבְרִיָּה וַעֲבָדְךָ שֵׁשׁ שָׁנִים וּבַשָּׁנָה הַשְּׁבִיעִת תְּשַׁלְּחֶנּוּ חפְשִׁי מֵעִמָּךְ׃}
}
{
	\eDtn{
		“When your Hebrew brother or sister will be sold to you,
		then he shall work for you six years, and in the seventh
		year you shall let him go liberated from you.
	}
}
{
%	\footnote{
%		let him go. This is the same term that was used for Pharaoh’s release of
%		the Israelite slaves: “Let my people go.” The Israelite should do no
%		less than Pharaoh! The Israelite should let his slave go, without the
%		coercion that Pharaoh needed, and after no more than six years.
%	}
}

\Verse{13}
{
	\hDtn{וְכִי תְשַׁלְּחֶנּוּ חפְשִׁי מֵעִמָּךְ לֹא תְשַׁלְּחֶנּוּ רֵיקָם׃}
}
{
	\eDtn{
		And when you let him go liberated from you, you shall not
		let him go empty-handed.
	}
}
{
}

\Verse{14}
{
	\hDtn{הַעֲנֵיק תַּעֲנִיק לוֹ מִצֹּאנְךָ וּמִגּרְנְךָ וּמִיִּקְבֶךָ אֲשֶׁר בֵּרַכְךָ יְהֹוָה אֱלֹהֶיךָ תִּתֶּן לוֹ׃}
}
{
	\eDtn{
		You shall provide him from your flock and from your
		threshing floor and from your winepress; as YHWH, your
		God, has blessed you, you shall give to him.
	}
}
{
}

\Verse{15}
{
	\hDtn{וְזָכַרְתָּ כִּי עֶבֶד הָיִיתָ בְּאֶרֶץ מִצְרַיִם וַיִּפְדְּךָ יְהֹוָה אֱלֹהֶיךָ עַל כֵּן אָנֹכִי מְצַוְּךָ אֶת הַדָּבָר הַזֶּה הַיּוֹם׃}
}
{
	\eDtn{
		And you shall remember that you were a slave in the land
		of Egypt, and YHWH, your God, redeemed you. On account of
		this, I command you this thing today.
	}
}
{
}

\Verse{16}
{
	\hDtn{וְהָיָה כִּי יֹאמַר אֵלֶיךָ לֹא אֵצֵא מֵעִמָּךְ כִּי אֲהֵבְךָ וְאֶת בֵּיתֶךָ כִּי טוֹב לוֹ עִמָּךְ׃}
}
{
	\eDtn{
		“And it will be, if he’ll say to you, ‘I won’t go out from
		you’ because he loves you and your house, because it is
		good for him with you,
	}
}
{
}

\Verse{17}
{
	\hDtn{וְלָקַחְתָּ אֶת הַמַּרְצֵעַ וְנָתַתָּה בְאזְנוֹ וּבַדֶּלֶת וְהָיָה לְךָ עֶבֶד עוֹלָם וְאַף לַאֲמָתְךָ תַּעֲשֶׂה כֵּן׃}
}
{
	\eDtn{
		then you shall take an awl and put it in his ear and in
		the door, and he shall be a slave forever to you. And you
		shall also do this to your maid.
	}
}
{
}

\Verse{18}
{
	\hDtn{לֹא יִקְשֶׁה בְעֵינֶךָ בְּשַׁלֵּחֲךָ אֹתוֹ חפְשִׁי מֵעִמָּךְ כִּי מִשְׁנֶה שְׂכַר שָׂכִיר עֲבָדְךָ שֵׁשׁ שָׁנִים וּבֵרַכְךָ יְהֹוָה אֱלֹהֶיךָ בְּכֹל אֲשֶׁר תַּעֲשֶׂה׃}
}
{
	\eDtn{
		It shall not be hard in your eyes when you let him go
		liberated from you, because he served you six years at
		twice the value of an employee, and YHWH, your God, will
		bless you in everything that you’ll do.
	}
}
{
}

\Verse{19}
{
	\hDtn{כּל הַבְּכוֹר אֲשֶׁר יִוָּלֵד בִּבְקָרְךָ וּבְצֹאנְךָ הַזָּכָר תַּקְדִּישׁ לַיהֹוָה אֱלֹהֶיךָ לֹא תַעֲבֹד בִּבְכֹר שׁוֹרֶךָ וְלֹא תָגֹז בְּכוֹר צֹאנֶךָ׃}
}
{
	\eDtn{
		“Every male firstborn that will be born in your herd and
		in your flock you shall consecrate to YHWH, your God. You
		shall not work with your ox’s firstborn, and you shall not
		shear your sheep’s firstborn.
	}
}
{
}

\Verse{20}
{
	\hDtn{לִפְנֵי יְהֹוָה אֱלֹהֶיךָ תֹאכְלֶנּוּ שָׁנָה בְשָׁנָה בַּמָּקוֹם אֲשֶׁר יִבְחַר יְהֹוָה אַתָּה וּבֵיתֶךָ׃}
}
{
	\eDtn{
		You shall eat it in front of YHWH, your God, year by year
		in the place that YHWH will choose—you and your household.
	}
}
{
}

\Verse{21}
{
	\hDtn{וְכִי יִהְיֶה בוֹ מוּם פִּסֵּחַ אוֹ עִוֵּר כֹּל מוּם רָע לֹא תִזְבָּחֶנּוּ לַיהֹוָה אֱלֹהֶיךָ׃}
}
{
	\eDtn{
		And if there will be an injury in it—crippled or blind,
		any bad injury—you shall not sacrifice it to YHWH, your
		God.
	}
}
{
}

\Verse{22}
{
	\hDtn{בִּשְׁעָרֶיךָ תֹּאכְלֶנּוּ הַטָּמֵא וְהַטָּהוֹר יַחְדָּו כַּצְּבִי וְכָאַיָּל׃}
}
{
	\eDtn{
		You shall eat it within your gates—the impure and the pure
		together, like a gazelle and like a deer.
	}
}
{
}

\Verse{23}
{
	\hDtn{רַק אֶת דָּמוֹ לֹא תֹאכֵל עַל הָאָרֶץ תִּשְׁפְּכֶנּוּ כַּמָּיִם׃}
}
{
	\eDtn{
		Only: you shall not eat its blood. You shall spill it like
		water on the ground.
	}
}
{
}
